\chapter{Beskrivende statistik}
\label{chap:ch6}

I softwareudvikling arbejder vi konstant med data: ydelsesmålinger, mønstre i brugeradfærd, systemers svartider og mål for kodekompleksitet. For at forstå disse oplysninger og uddrage meningsfuld indsigt har vi brug for værktøjer til at opsummere, organisere og visualisere data. Beskrivende statistik giver os disse grundlæggende værktøjer, så vi kan forstå vores datas egenskaber uden at drage slutninger om større populationer.

Forestil dig et softwareudviklingsteam, der analyserer ydelsen af sin webapplikation. Teamet indsamler svartider for 1000 brugerforespørgsler og finder værdier fra 50 millisekunder til 3.2 sekunder. Uden en ordentlig analyse er disse rådata overvældende og siger ikke meget. Ved at anvende beskrivende statistik kan teamet derimod fastslå, at den gennemsnitlige svartid er 450 millisekunder, at 75\% af forespørgslerne er afsluttet inden for 600 millisekunder, og at der er nogle få usædvanligt langsomme forespørgsler, som kan tyde på ydelsesproblemer.

Dette kapitel introducerer de grundlæggende begreber i beskrivende statistik med fokus på mål, der hjælper os med at forstå centrale tendenser, variation og fordelingsmønstre i vores data. Vi undersøger, hvordan disse teknikker anvendes specifikt i softwareudvikling, fra analyse af algoritmers ydelse til forståelse af mønstre i brugeradfærd.

\section{Introduktion til beskrivende statistik} \label{sec:ch6sec1}

Nogle gange begynder statistisk arbejde med eksisterende data, fx nedbørsmålinger, ledighedsprocenter eller BNP-tal, som vi derefter opsummerer og analyserer. I andre situationer skal data frembringes gennem et eksperiment eller en undersøgelse. Hvis en underviser fx vil sammenligne to undervisningsmetoder i et introducerende programmeringskursus, kan underviseren tilfældigt dele de studerende i to grupper, anvende en forskellig metode i hver gruppe og derefter sammenligne testresultaterne. Den tilfældige fordeling er afgørende: Den sikrer, at forskelle mellem grupperne ikke skyldes forhold, der fandtes i forvejen, fx evner, men derimod selve undervisningsmetoden.  

Når data er indsamlet, bliver de opsummeret og visualiseret -- for eksempel ved at angive det gennemsnitlige resultat for hver gruppe. Denne proces er kernen i \textbf{beskrivende statistik}: at organisere, præsentere og beskrive data på en meningsfuld måde.  

\subsection*{Population og stikprøve}

I statistisk analyse skelner vi mellem to centrale begreber:

\begin{definition}[Population]
En \textbf{population} er den fuldstændige mængde af alle mulige observationer eller målinger, der er af interesse i en bestemt undersøgelse. I softwareudvikling kan det være alle mulige køretider for en algoritme, alle brugersessioner på et websted eller alle kodelinjer i et projekt.
\end{definition}

\begin{definition}[Stikprøve]
En \textbf{stikprøve} er en delmængde af populationen, som vi faktisk observerer eller måler. På grund af praktiske begrænsninger arbejder vi ofte med stikprøver i stedet for hele populationer.
\end{definition}

I eksemplet med undervisningsmetoderne udgør de studerende i klasselokalet kun en \emph{stikprøve} af den større \emph{population} af potentielle studerende. Beskrivende statistik hjælper os med at forstå selve stikprøven, mens inferentielle metoder senere gør det muligt at bruge stikprøven til at drage konklusioner om den bredere population. Inferentiel statistik ligger uden for denne bogs rammer. I dette kapitel fokuserer vi på beskrivende statistik.


\begin{example} Analyse af softwareydelse

Se på en webapplikation, der betjener millioner af brugere dagligt. \textbf{Populationen} vil være svartiderne for alle mulige brugerforespørgsler. På grund af begrænsninger i beregnings- og lagerkapacitet indsamler vi dog måske kun svartider for en \textbf{stikprøve} på 10\,000 forespørgsler om dagen. Stikprøven bør være repræsentativ for populationen, hvis vi skal kunne drage meningsfulde konklusioner.
\end{example}

\subsection*{Datatyper}

Data kan inddeles i forskellige typer, som hver kræver forskellige statistiske tilgange:

\begin{definition}[Kvalitative data]
\textbf{Kvalitative data} (også kaldet kategoriske data) består af ikke-numeriske oplysninger, der kan inddeles i kategorier. Eksempler er de programmeringssprog, der bruges i et projekt, brugertilfredshed (tilfreds/neutral/utilfreds) eller fejls alvorlighedsgrad (kritisk/høj/middel/lav).
\end{definition}

Den væsentlige forskel mellem \textbf{kvalitative} og \textbf{kvantitative} data ligger i, hvad de repræsenterer, og hvordan de analyseres. Kvalitative data er kategorier eller egenskaber, der ikke meningsfuldt kan udtrykkes som tal. Sådanne data bruges til at klassificere eller mærke elementer, fx programmeringssprog eller fejls alvorlighedsgrad. Kvantitative data består derimod af numeriske værdier, der repræsenterer antal eller målinger, fx svartider eller antal brugere. Fordi disse værdier er numeriske, kan man udføre matematiske operationer og statistisk analyse på dem. Det er grundlæggende at kunne skelne mellem kvalitative og kvantitative data, da det styrer valget af statistiske metoder og af, hvilke typer visualiseringer der er egnede til analysen.

\begin{definition}[Kvantitative data]
\textbf{Kvantitative data} består af numeriske målinger, der kan ordnes og underkastes matematiske operationer. De kan yderligere opdeles i:
\begin{itemize}
    \item \textbf{Diskrete}: Tællelige værdier (antal fejl, kodelinjer, brugersessioner)
    \item \textbf{Kontinuerte}: Målbare værdier, der kan antage enhver værdi inden for et interval (svartider, hukommelsesforbrug, CPU-udnyttelse)
\end{itemize}
\end{definition}

\begin{example}
\begin{itemize}
    \item \textbf{Kvalitative}: Programmeringssprog (Python, Java, C++), driftsmiljø (udvikling, staging, produktion)
    \item \textbf{Kvantitative, diskrete}: Antal commits pr. dag, kodelinjer, antal testtilfælde
    \item \textbf{Kvantitative, kontinuerte}: Svartid i millisekunder, hukommelsesforbrug i MB, CPU-udnyttelse i procent
\end{itemize}
\end{example}

\subsection*{Outliers og ekstreme data}

I ethvert datasæt har de fleste værdier en tendens til at samle sig i et centralt område, men af og til møder vi værdier, der er meget større eller mindre end resten. De kaldes \textbf{outliers} eller \textbf{ekstremværdier}. Outliers kan opstå af mange grunde: målefejl, usædvanlige, men gyldige hændelser eller naturlig variation i data.

Hvis de fleste webforespørgsler fx er afsluttet på under 400 millisekunder, men én forespørgsel tager 1200 millisekunder, er værdien 1200~ms en outlier. Outliers kan have stor indflydelse på opsummerende statistikker som gennemsnittet, så data ser ud til at variere mere, eller så gennemsnittet flyttes væk fra det sted, hvor de fleste værdier ligger.

Det er vigtigt at holde øje med outliers, når man analyserer data, da de kan være tegn på interessante fænomener (fx en sjælden flaskehals i ydelsen), indtastningsfejl eller et behov for nærmere undersøgelse. Der findes formelle metoder til at finde outliers, som vi gennemgår senere, men ofte er et simpelt plot eller et blik på de sorterede data nok til at få øje på værdier, der »skiller sig ud« fra resten.

I praksis er det afgørende at forstå sammenhængen: Nogle gange er outliers fejl, der skal rettes eller fjernes, men andre gange er de den vigtigste del af historien.

\section{Beskrivelse af datasæt}
Når man præsenterer numeriske resultater, er klarhed og kortfattethed afgørende. Tabeller og grafer er særligt effektive, fordi de lader læseren hurtigt få fat i et datasæts vigtigste egenskaber. Visuelle opsummeringer og tabeller fremhæver ofte vigtige forhold som den samlede variationsbredde, hvor tæt værdierne er samlet, og om data ser symmetriske eller skæve ud.

\subsection*{Hyppighedstabeller og grafer} \label{sub:frequency_tables_and_graphs}

Det første skridt i en dataanalyse er ofte at organisere rådata i en hyppighedsfordeling, som viser, hvor ofte hver værdi (eller hvert interval af værdier) forekommer.

\begin{definition}[Hyppighedsfordeling]
En \textbf{hyppighedsfordeling} er en tabel, der viser hyppigheden (antallet) af hver værdi eller klasse af værdier i et datasæt. Den kan angives som:
\begin{itemize}
    \item \textbf{Absolut hyppighed}: Det faktiske antal forekomster
    \item \textbf{Relativ hyppighed}: Andelen af det samlede antal observationer
    \item \textbf{Kumuleret hyppighed}: Den løbende sum af hyppighederne
\end{itemize}
\end{definition}

\begin{example}
Et datasæt med et forholdsvis lille antal forskellige værdier kan bekvemt præsenteres i en hyppighedstabel. Se fx på et datasæt med de månedlige startlønninger (afrundet til nærmeste tusind danske kroner) for 42 nyuddannede med en bachelorgrad i softwareudvikling.

Rådata er som følger:
\[
\begin{aligned}
&37, 37, 37, 37, 38, 39, 39, 39, 40, 40, 40, 40, 40, 41, 41, 41, 41, 41, 41, 41, 41, \\
&42, 42, 42, 42, 42, 42, 42, 42, 42, 42, 44, 44, 44, 44, 44, 46, 46, 47, 47, 47, 50
\end{aligned}
\]

Ud fra disse data kan vi opstille en hyppighedstabel. \autoref{tab:salary_freq} fortæller os blandt andet, at den laveste startløn på 37\,000 DKK blev givet til fire af de nyuddannede, mens den højeste løn på 50\,000 DKK kun blev givet til én. Den hyppigste startløn var 42\,000 DKK, som 10 af de nyuddannede fik.

\begin{table}[htbp]
\centering
\renewcommand{\arraystretch}{1.2}
{\arrayrulecolor{mseViaBlue}
\begin{tabular}{cc}
\toprule
\textbf{Løn (\(\times 1000\) DKK)} & \textbf{Hyppighed} \\
\midrule
37 & 4 \\
38 & 1 \\
39 & 3 \\
40 & 5 \\
41 & 8 \\
42 & 10 \\
44 & 5 \\
46 & 2 \\
47 & 3 \\
50 & 1 \\
\bottomrule
\end{tabular}
}
\caption{Hyppighedsfordeling af månedlige startlønninger}
\label{tab:salary_freq}
\end{table}
\end{example}

Data fra en hyppighedstabel kan afbildes grafisk i et pindediagram, der afsætter de
forskellige dataværdier på den vandrette akse og angiver deres hyppigheder med højden af
lodrette linjer. Et pindediagram over dataene i \autoref{tab:salary_freq} er vist i \autoref{fig:salary_line_graph}.

\begin{figure}[htbp]
\centering
\begin{tikzpicture}
\begin{axis}[
    title={Pindediagram over startlønninger},
    xlabel={Startløn (\(\times 1000\) DKK)},
    ylabel={Hyppighed},
    ymin=0,
    ymax=12,
    xmin=36,
    xmax=51,
    grid=major,
    grid style={dashed,gray!30},
    xtick={37,38,...,50},
    xticklabel style={rotate=45, anchor=east, yshift=-5pt},
    width=0.9\linewidth,
    height=7cm,
    title style={font=\large\bfseries}
]
\addplot+[ycomb, mseViaBlue, thick, mark=*, mark options={fill=mseViaBlue,draw=mseViaBlue}] coordinates {
    (37, 4)
    (38, 1)
    (39, 3)
    (40, 5)
    (41, 8)
    (42, 10)
    (44, 5)
    (46, 2)
    (47, 3)
    (50, 1)
};
\end{axis}
\end{tikzpicture}
\caption{Et pindediagram, der viser hyppigheden af de forskellige startlønninger.}
\label{fig:salary_line_graph}
\end{figure}

Når linjerne i et pindediagram gøres tykkere, kaldes grafen et søjlediagram. \autoref{fig:salary_bar_graph} viser et søjlediagram over de samme lønninger.

\newpage   

\begin{figure}[htbp]
\centering
\begin{tikzpicture}
\begin{axis}[
    title={Søjlediagram over startlønninger},
    xlabel={Startløn (\(\times 1000\) DKK)},
    ylabel={Hyppighed},
    ymin=0,
    ymax=12,
    xmin=36,
    xmax=51,
    grid=major,
    grid style={dashed,gray!30},
    xtick={37,38,...,50},
    xticklabel style={rotate=45, anchor=east, yshift=-5pt},
    width=0.9\linewidth,
    height=7cm,
    title style={font=\large\bfseries},
    ybar,
    bar width=0.6cm
]
\addplot[fill=mseViaBlue, draw=mseViaBlue!80!black] coordinates {
    (37, 4)
    (38, 1)
    (39, 3)
    (40, 5)
    (41, 8)
    (42, 10)
    (43, 0)
    (44, 5)
    (46, 2)
    (47, 3)
    (50, 1)
};
\end{axis}
\end{tikzpicture}
\caption{Et søjlediagram over lønningerne.}
\label{fig:salary_bar_graph}
\end{figure}

En anden type graf, der bruges til at afbilde en hyppighedstabel, er hyppighedspolygonen, som
afsætter hyppighederne af de forskellige dataværdier på den lodrette akse og derefter forbinder de
afsatte punkter med rette linjestykker. \autoref{fig:salary_frequency_polygon} viser en hyppighedspolygon for de samme lønninger.

\begin{figure}[htbp]
\centering
\begin{tikzpicture}
\begin{axis}[
    title={Hyppighedspolygon for startlønninger},
    xlabel={Startløn (\(\times 1000\) DKK)},
    ylabel={Hyppighed},
    ymin=0,
    ymax=12,
    xmin=36,
    xmax=51,
    grid=major,
    grid style={dashed,gray!30},
    xtick={37,38,...,50},
    xticklabel style={rotate=45, anchor=east, yshift=-5pt},
    width=0.9\linewidth,
    height=7cm,
    title style={font=\large\bfseries}
]
\addplot[color=mseViaBlue, thick, mark=*, mark options={fill=mseViaBlue}] coordinates {
    (37, 4)
    (38, 1)
    (39, 3)
    (40, 5)
    (41, 8)
    (42, 10)
    (43, 0)
    (44, 5)
    (45, 0)
    (46, 2)
    (47, 3)
    (48, 0)
    (49, 0)
    (50, 1)
};
\end{axis}
\end{tikzpicture}
\caption{En hyppighedspolygon for lønningerne, der forbinder punkterne fra hyppighedstabellen.}
\label{fig:salary_frequency_polygon}
\end{figure}

Nogle gange foretrækker vi den relative hyppighed frem for den absolutte hyppighed. Det er særligt nyttigt ved store datasæt, og når man vil sammenligne værdier fra datasæt af forskellig størrelse. En udvidelse af den relative hyppighed er den kumulerede relative hyppighed, som er den løbende sum af de relative hyppigheder. \autoref{tab:salary_full_freq} viser hyppigheden, den relative hyppighed og den kumulerede relative hyppighed for lønningerne.

\begin{table}[htbp]
\centering
\renewcommand{\arraystretch}{1.3}
{\arrayrulecolor{mseViaBlue}
\begin{tabular}{cccc}
\toprule
\textbf{Løn} & \textbf{Hyppighed} & \textbf{Relativ hyppighed} & \textbf{Kumuleret rel. hyppighed} \\
\midrule
37 & 4 & $4/42 \approx 0.095$ & $4/42 \approx 0.095$ \\
38 & 1 & $1/42 \approx 0.024$ & $5/42 \approx 0.119$ \\
39 & 3 & $3/42 \approx 0.071$ & $8/42 \approx 0.190$ \\
40 & 5 & $5/42 \approx 0.119$ & $13/42 \approx 0.310$ \\
41 & 8 & $8/42 \approx 0.190$ & $21/42 = 0.500$ \\
42 & 10 & $10/42 \approx 0.238$ & $31/42 \approx 0.738$ \\
44 & 5 & $5/42 \approx 0.119$ & $36/42 \approx 0.857$ \\
46 & 2 & $2/42 \approx 0.048$ & $38/42 \approx 0.905$ \\
47 & 3 & $3/42 \approx 0.071$ & $41/42 \approx 0.976$ \\
50 & 1 & $1/42 \approx 0.024$ & $42/42 = 1.000$ \\
\bottomrule
\end{tabular}
}
\caption{Hyppighed, relativ hyppighed og kumuleret relativ hyppighed af månedlige startlønninger.}
\label{tab:salary_full_freq}
\end{table}

\newpage

\subsection*{Grupperede data og histogrammer}

Som vi så i \autoref{sub:frequency_tables_and_graphs}, er pindediagrammer og søjlediagrammer nyttige til at vise dataværdiers hyppigheder. Når et datasæt indeholder mange forskellige værdier, som pærernes levetider i \autoref{tab:lamp_life_raw}, bliver det upraktisk at lave en hyppighedstabel for hver enkelt værdi. En mere effektiv fremgangsmåde er at gruppere data i et antal \textbf{klasseintervaller}, som også ofte kaldes \textbf{intervaller} (på engelsk \emph{bins}).

\begin{table}[h!]
\centering
\small % Use a smaller font for the large table
\renewcommand{\arraystretch}{1.2}
{\arrayrulecolor{mseViaBlue}
\begin{tabular}{rrrrrrrrrr}
\toprule
\multicolumn{10}{c}{\textbf{Levetider (i timer)}} \\
\midrule
1067 & 919 & 1196 & 785 & 1126 & 936 & 918 & 1156 & 920 & 948 \\
855 & 1092 & 1162 & 1170 & 929 & 950 & 905 & 972 & 1035 & 1045 \\
1157 & 1195 & 1195 & 1340 & 1122 & 938 & 970 & 1237 & 956 & 1102 \\
1022 & 978 & 832 & 1009 & 1157 & 1151 & 1009 & 765 & 958 & 902 \\
923 & 1333 & 811 & 1217 & 1085 & 896 & 958 & 1311 & 1037 & 702 \\
521 & 933 & 928 & 1153 & 946 & 858 & 1071 & 1069 & 830 & 1063 \\
930 & 807 & 954 & 1063 & 1002 & 909 & 1077 & 1021 & 1062 & 1157 \\
999 & 932 & 1035 & 944 & 1049 & 940 & 1122 & 1115 & 833 & 1320 \\
901 & 1324 & 818 & 1250 & 1203 & 1078 & 890 & 1303 & 1011 & 1102 \\
996 & 780 & 900 & 1106 & 704 & 621 & 854 & 1178 & 1138 & 951 \\
1187 & 1067 & 1118 & 1037 & 958 & 760 & 1101 & 949 & 992 & 966 \\
824 & 653 & 980 & 935 & 878 & 934 & 910 & 1058 & 730 & 980 \\
844 & 814 & 1103 & 1000 & 788 & 1143 & 935 & 1069 & 1170 & 1067 \\
1037 & 1151 & 863 & 990 & 1035 & 1112 & 931 & 970 & 932 & 904 \\
1026 & 1147 & 883 & 867 & 990 & 1258 & 1192 & 922 & 1150 & 1091 \\
1039 & 1083 & 1040 & 1289 & 699 & 1083 & 880 & 1029 & 658 & 912 \\
1023 & 984 & 856 & 924 & 801 & 1122 & 1292 & 1116 & 880 & 1173 \\
1134 & 932 & 938 & 1078 & 1180 & 1106 & 1184 & 954 & 824 & 529 \\
998 & 996 & 1133 & 765 & 775 & 1105 & 1081 & 1171 & 705 & 1425 \\
610 & 916 & 1001 & 895 & 709 & 860 & 1110 & 1149 & 972 & 1002 \\
\bottomrule
\end{tabular}
}
\caption{Levetid i timer for 200 glødepærer.}
\label{tab:lamp_life_raw}
\end{table}

Valget af, hvor mange intervaller man skal bruge, er en afvejning. Med for få intervaller risikerer vi at miste vigtig information, fordi for mange forskellige værdier bliver grupperet sammen. Med for mange intervaller kan hyppighederne i hvert interval blive så små, at de ikke afslører noget tydeligt mønster i datas fordeling.

Det optimale antal intervaller er et subjektivt valg, der afhænger af datasættet, men det er typisk at vælge 5 til 10 intervaller. Det er også almindeligt, men ikke påkrævet, at bruge lige brede intervaller, da det ofte gør fortolkningen lettere.

Et intervals endepunkter kaldes dets \textbf{klassegrænser}. Vi følger den almindelige \textbf{konvention om inklusion af venstre endepunkt}, som fastlægger, at et interval indeholder sin venstre grænse, men ikke sin højre. Intervallet 700--800 indeholder for eksempel enhver værdi \(x\), hvor \(700 \le x < 800\).

Med denne fremgangsmåde er rådata fra \autoref{tab:lamp_life_raw} opsummeret i \autoref{tab:lamp_life_freq} med intervaller af længden 100, startende ved 500.

\begin{table}[htbp]
\centering
\renewcommand{\arraystretch}{1.3}
{\arrayrulecolor{mseViaBlue}
\begin{tabular}{cccc}
\toprule
\textbf{Klasseinterval} & \textbf{Hyppighed} & \textbf{Relativ hyppighed} & \textbf{Kumuleret rel. hyppighed} \\
\midrule
500--600   & 2   & 0.010 & 0.010 \\
600--700   & 5   & 0.025 & 0.035 \\
700--800   & 12  & 0.060 & 0.095 \\
800--900   & 25  & 0.125 & 0.220 \\
900--1000  & 58  & 0.290 & 0.510 \\
1000--1100 & 41  & 0.205 & 0.715 \\
1100--1200 & 43  & 0.215 & 0.930 \\
1200--1300 & 7   & 0.035 & 0.965 \\
1300--1400 & 6   & 0.030 & 0.995 \\
1400--1500 & 1   & 0.005 & 1.000 \\
\bottomrule
\end{tabular}
}
\caption{Hyppighedsfordeling for levetiderne af 200 glødepærer.}
\label{tab:lamp_life_freq}
\end{table}

Et søjlediagram, der viser hyppighederne af data grupperet i intervaller, kaldes et \textbf{histogram}. Et vigtigt kendetegn ved et histogram er, at søjlerne står direkte op ad hinanden, hvilket afspejler, at intervallerne er sammenhængende. Den lodrette akse kan angive enten den absolutte hyppighed (antallet af værdier i hvert interval) eller den relative hyppighed. Det første kaldes et \textbf{hyppighedshistogram}, det sidste et \textbf{relativt hyppighedshistogram}. \autoref{fig:lamp_life_histogram} viser et hyppighedshistogram for pærernes levetider.

\newpage

\begin{figure}[htbp]
\centering
\begin{tikzpicture}
\begin{axis}[
    title={Hyppighedshistogram over pærernes levetider},
    xlabel={Levetid (i timer)},
    ylabel={Hyppighed},
    ymin=0,
    width=0.9\linewidth,
    height=7cm,
    ybar interval,
    xticklabel={\pgfmathprintnumber\tick--\pgfmathprintnumber\nexttick},
    xticklabel style={rotate=45, anchor=east, yshift=-5pt},
    grid=major,
    grid style={dashed,gray!30},
    title style={font=\large\bfseries}
]
\addplot[fill=mseViaBlue, draw=mseViaBlue!80!black] coordinates {
    (500, 2)
    (600, 5)
    (700, 12)
    (800, 25)
    (900, 58)
    (1000, 41)
    (1100, 43)
    (1200, 7)
    (1300, 6)
    (1400, 1)
    (1500, 0)
};
\end{axis}
\end{tikzpicture}
\caption{Et hyppighedshistogram for pærernes levetider. De tilstødende søjler viser, hvordan data fordeler sig på sammenhængende intervaller.}
\label{fig:lamp_life_histogram}
\end{figure}

Endelig er vi nogle gange interesserede i at afbilde den kumulerede hyppighed eller, oftere, den kumulerede relative hyppighed. En sådan graf, der ofte kaldes en \textbf{sumkurve} (på engelsk \emph{ogive}), viser antallet eller andelen af datapunkter, der ligger på eller under en bestemt værdi.

\begin{figure}[h!]
\centering
\begin{tikzpicture}
\begin{axis}[
    title={Kumuleret relativ hyppighed af pærernes levetider (sumkurve)},
    xlabel={Levetid (i timer)},
    ylabel={Kumuleret relativ hyppighed},
    ymin=0, ymax=1.05,
    xmin=400, xmax=1500,
    width=0.9\linewidth,
    height=7cm,
    grid=major,
    grid style={dashed,gray!30},
    title style={font=\large\bfseries}
]
\addplot[color=mseViaBlue, thick, mark=*, mark options={fill=mseViaBlue}] coordinates {
    (500, 0)
    (600, 0.010)
    (700, 0.035)
    (800, 0.095)
    (900, 0.220)
    (1000, 0.510)
    (1100, 0.715)
    (1200, 0.930)
    (1300, 0.965)
    (1400, 0.995)
    (1500, 1.000)
};
\end{axis}
\end{tikzpicture}
\caption{En sumkurve, der viser den kumulerede relative hyppighed for pærernes levetider.}
\label{fig:lamp_life_ogive}
\end{figure}

På en sumkurve repræsenterer hvert punkt på den vandrette akse en dataværdi, mens det tilsvarende punkt på den lodrette akse viser andelen af data, der er mindre end eller lig med denne værdi. \autoref{fig:lamp_life_ogive} viser fx den kumulerede relative hyppighed for pærernes levetider. Ud fra grafen kan vi hurtigt se, at omkring 50\% af pærerne havde en levetid på under 1000 timer, at cirka 95\% havde en levetid på under 1200 timer (den faktiske værdi er 93\%), og at 100\% af pærerne var gået i stykker efter 1500 timer.

\section{Opsummering af data}
Vi har set, hvordan man kan opsummere datafordelinger visuelt og numerisk. Nu fokuserer vi på to grundlæggende aspekter af beskrivelsen af et datasæt: dets \emph{centrum} (eller typiske værdi) og dets \emph{spredning} (eller variation).

\subsection*{Centralmål}
\textbf{Gennemsnittet} (også kaldet middelværdien) er det mest anvendte centralmål. \autoref{def:arithmetic_mean} viser, hvordan man beregner et datasæts gennemsnit, og \autoref{ex:mean_response_time} giver et eksempel.

\begin{definition}[Aritmetisk gennemsnit] \label{def:arithmetic_mean}

For et datasæt med $n$ værdier $x_1, x_2, \ldots, x_n$ er det \textbf{aritmetiske gennemsnit} defineret som:

$$
\bar{x} = \frac{1}{n}\sum_{i=1}^{n} x_i = \frac{x_1 + x_2 + \cdots + x_n}{n}
$$

\end{definition}


\begin{example} Analyse af svartider \label{ex:mean_response_time}

En webapplikation har registreret følgende svartider (i millisekunder) for 10 forespørgsler:

\[
\{120, 150, 180, 200, 220, 250, 280, 300, 350, 400\}
\]

Den gennemsnitlige svartid er:
\[
\bar{x} = \frac{120 + 150 + 180 + 200 + 220 + 250 + 280 + 300 + 350 + 400}{10} = \frac{2450}{10} = 245 \text{ millisekunder}
\]
\end{example}

Gennemsnittet inddrager alle værdier i et datasæt, hvilket gør det til et stærkt centralmål. Det er mest velegnet, når dataværdierne er af nogenlunde samme størrelsesorden, og når fordelingen er rimeligt symmetrisk uden ekstreme outliers. Under disse betingelser giver gennemsnittet en pålidelig opsummering af datasættet.

I analysen af webydelse giver den gennemsnitlige svartid fx et samlet billede af, hvor hurtigt et system i gennemsnit svarer. Hvis nogle få svartider er meget langsommere end resten, kan gennemsnittet dog give et misvisende indtryk af den typiske brugeroplevelse.

Ud over sin rolle som opsummerende statistik er gennemsnittet også grundlaget for mange statistiske metoder. Det er især centralt i definitionerne af varians og standardafvigelse, som måler datas spredning omkring gennemsnittet.

\begin{remark}
    Som nævnt tidligere, og som vi vender tilbage til, er gennemsnittet meget følsomt over for ekstremværdier eller outliers. Selv en enkelt usædvanligt stor eller lille observation kan flytte gennemsnittet betydeligt, så det i mindre grad afspejler datasættets typiske værdi.
    Det er derfor vigtigt at se på datas fordeling, før man udelukkende bruger gennemsnittet til at beskrive et datasæts centrum.
\end{remark}

\textbf{Medianen} er den midterste værdi, når data er ordnet i stigende rækkefølge. Den er mindre følsom over for outliers end gennemsnittet.

\begin{definition}[Median]
    For et datasæt med $n$ værdier ordnet i stigende rækkefølge er medianen defineret som:
    \[
    \text{median} =
    \begin{cases}
        x_{\frac{n+1}{2}} & \text{hvis } n \text{ er ulige} \\
        \frac{x_{\frac{n}{2}} + x_{\frac{n}{2}+1}}{2} & \text{hvis } n \text{ er lige}
    \end{cases}
    \]
\end{definition}

\begin{example} Median svartid

Vi bruger de samme svartider som i det foregående eksempel:
\[
\{120, 150, 180, 200, 220, 250, 280, 300, 350, 400\}
\]

Da $n = 10$ (lige), er medianen gennemsnittet af den 5. og 6. værdi:
\[
\text{Median} = \frac{220 + 250}{2} = 235 \text{ millisekunder}
\]
\end{example}

\begin{example}

Se på de samme data med en ekstremværdi:
\[
\{120, 150, 180, 200, 220, 250, 280, 300, 350, 2000\}
\]

Gennemsnittet bliver: $\bar{x} = \frac{4050}{10} = 405$ millisekunder (stærkt påvirket af outlieren)

Medianen forbliver: $\frac{220 + 250}{2} = 235$ millisekunder (upåvirket af outlieren)
\end{example}

Fordi medianen deler et datasæt i to lige store halvdele, er den et robust centralmål, især når data indeholder outliers eller er skæve. I modsætning til gennemsnittet, som kan trækkes mod ekstremværdier, er medianen stort set upåvirket og giver i sådanne situationer et mere pålideligt billede af den »typiske« værdi.

\textbf{Typetallet} er den værdi, der forekommer oftest i et datasæt.

\begin{definition}[Typetal]
\textbf{Typetallet} er den værdi, der forekommer hyppigst i et datasæt. Et datasæt kan have:
\begin{itemize}
    \item \textbf{Intet typetal}: Hvis alle værdier forekommer lige hyppigt
    \item \textbf{Ét typetal (unimodalt)}: Hvis én værdi forekommer hyppigst
    \item \textbf{Flere typetal (multimodalt)}: Hvis to eller flere værdier deler den højeste hyppighed
\end{itemize}
\end{definition}

Når et datasæt har to typetal, kaldes det nogle gange et \textbf{bimodalt} datasæt.

\newpage

\begin{example}

Et softwareprojekt har følgende alvorlighedsgrader for fejl:
\[
\{\text{Høj}, \text{Middel}, \text{Lav}, \text{Høj}, \text{Kritisk}, \text{Middel}, \text{Høj}, \text{Lav}, \text{Høj}, \text{Middel}\}
\]

Optælling af hyppigheder:
\begin{itemize}
    \item Høj: 4 forekomster
    \item Middel: 3 forekomster  
    \item Lav: 2 forekomster
    \item Kritisk: 1 forekomst
\end{itemize}

Typetallet er »Høj«, da denne værdi forekommer hyppigst.
\end{example}

Typetallet er især nyttigt ved analyse af kategoriske eller diskrete data, hvor det ikke altid giver mening at beregne et gennemsnit eller en median. I softwareudvikling kan typetallet fx hjælpe med at finde den fejlkode, som et API oftest returnerer, det mest anvendte programmeringssprog i en kodebase eller den hyppigst rapporterede fejltype i et issue-tracking-system.

\subsubsection*{Hvornår bruges hvilket mål?}

\begin{table}[htbp]
\centering
\renewcommand{\arraystretch}{1.4}
\begin{tabular}{|l|l|l|}
\hline
\textbf{Mål} & \textbf{Bedst til} & \textbf{Begrænsninger} \\
\hline
Gennemsnit & Kontinuerte data, symmetriske fordelinger & Følsomt over for outliers \\
\hline
Median & Data med outliers, skæve fordelinger & Mindre informativ ved symmetriske data \\
\hline
Typetal & Kategoriske data, diskrete data & Findes ikke altid eller er ikke entydigt \\
\hline
\end{tabular}
\caption{Sammenligning af centralmål}
\label{tab:central_tendency}
\end{table}

\subsection*{Spredningsmål}

Mens centralmål fortæller os om den typiske værdi, beskriver \textbf{spredningsmål}, hvor spredte data er. I softwareudvikling er det afgørende at forstå variation, når man skal vurdere konsistens og pålidelighed.

\textbf{Variationsbredden} er det simpleste spredningsmål.

\begin{definition}[Variationsbredde]
Et datasæts \textbf{variationsbredde} er forskellen mellem den største og den mindste værdi:
\[
\text{Variationsbredde} = x_{\max} - x_{\min}
\]
\end{definition}

\begin{example}

    For svartiderne: $\{120, 150, 180, 200, 220, 250, 280, 300, 350, 400\}$

\[
\text{Variationsbredde} = 400 - 120 = 280 \text{ millisekunder}
\]
\end{example}

\begin{remark}
Variationsbredden er følsom over for outliers og giver ingen information om, hvordan værdierne er fordelt mellem yderpunkterne.
\end{remark}

\textbf{Variansen} og \textbf{standardafvigelsen} er mere avancerede mål, der tager hensyn til alle datapunkter.

\begin{definition}[Stikprøvevarians og standardafvigelse]
For en stikprøve med $n$ værdier er \textbf{stikprøvevariansen}, betegnet $s^2$, defineret som:
$$
s^2 = \frac{1}{n-1}\sum_{i=1}^{n}(x_i - \bar{x})^2
$$
\textbf{Stikprøvestandardafvigelsen}, betegnet $s$, er kvadratroden af variansen:
$$
s = \sqrt{s^2} = \sqrt{\frac{1}{n-1}\sum_{i=1}^{n}(x_i - \bar{x})^2}
$$
\end{definition}

Mens gennemsnittet angiver et datasæts centrum, måler \textbf{variansen} og \textbf{standardafvigelsen} dets variation eller spredning. En lille standardafvigelse betyder, at datapunkterne ligger tæt samlet omkring gennemsnittet, mens en stor standardafvigelse betyder, at de er mere spredte.

Stikprøvevariansen er en skaleret sum af de kvadrerede afvigelser fra gennemsnittet med divisoren $n-1$. Kvadreringen af hver afvigelse tjener to formål: Den gør alle bidrag ikke-negative, og den giver større afvigelser større vægt. Beregningen efterlader dog variansen i kvadrerede enheder (fx millisekunder i anden), som ikke er intuitive at fortolke.

Det løser vi ved at bruge standardafvigelsen, som blot er kvadratroden af variansen. Dette afgørende skridt bringer målet tilbage til datas oprindelige enheder (fx millisekunder), så det kan fortolkes direkte. Kort sagt repræsenterer standardafvigelsen et datapunkts typiske afstand fra gennemsnittet.

\begin{example} Beregning af standardafvigelse

    For svartiderne: $\{120, 150, 180, 200, 220, 250, 280, 300, 350, 400\}$ med $\bar{x} = 245$:

\begin{align*}
s^2 &= \frac{1}{9}[(120-245)^2 + (150-245)^2 + \cdots + (400-245)^2] \\
&= \frac{1}{9}[(-125)^2 + (-95)^2 + \cdots + (155)^2] \\
&= \frac{1}{9}[15625 + 9025 + \cdots + 24025] \\
    &= \frac{1}{9} \times 70850 \approx 7872.22
\end{align*}

\[
s = \sqrt{7872.22} \approx 88.73 \text{ millisekunder}
\]
\end{example}

Følgende algebraiske identitet er ofte nyttig, når man beregner varians eller standardafvigelse i hånden:

Vi tager udgangspunkt i definitionen:

\[
s^2 = \frac{1}{n-1}\sum_{i=1}^n (x_i - \bar{x})^2
\]

Ved hjælp af identiteten

\[
\sum_{i=1}^n (x_i - \bar{x})^2 = \sum_{i=1}^n x_i^2 - \frac{\left(\sum_{i=1}^n x_i \right)^2}{n}
\]

kan vi omskrive stikprøvevariansen til

\[
s^2 = \frac{1}{n-1} \left( \sum_{i=1}^n x_i^2 - \frac{\left(\sum_{i=1}^n x_i\right)^2}{n} \right) =\frac{1}{n-1}\left(\sum_{i=1}^n x_i^2-n \bar{x}^2\right)
\]

Dette kaldes \textbf{regneformlen for varians}, da den gør det muligt at beregne $s^2$ direkte ud fra $\sum x_i$ og $\sum x_i^2$ (eller tilsvarende ud fra $\sum x_i^2$ og $\bar{x}$) uden at skulle beregne hver enkelt afvigelse fra gennemsnittet.


\textbf{Variationskoefficienten} gør det muligt at sammenligne variation på tværs af datasæt med forskellige skalaer.

\begin{definition}[Variationskoefficient]
\textbf{Variationskoefficienten} er forholdet mellem standardafvigelsen og gennemsnittet:
\[
CV = \frac{s}{\bar{x}} \times 100\%
\]
\end{definition}

\begin{example} Sammenligning af variation

Se på to systemer:
\begin{itemize}
    \item System A: Gennemsnitlig svartid = 100~ms, standardafvigelse = 20~ms
    \item System B: Gennemsnitlig svartid = 500~ms, standardafvigelse = 100~ms
\end{itemize}

\begin{align*}
CV_A &= \frac{20}{100} \times 100\% = 20\% \\
CV_B &= \frac{100}{500} \times 100\% = 20\%
\end{align*}

Begge systemer har samme relative variation (20\%), selv om deres absolutte standardafvigelser er forskellige.
\end{example}

\subsection*{Positionsmål}

For at kunne introducere et spredningsmål, der er mindre følsomt over for outliers, skal vi først forstå positionsmål, som beskriver et datapunkts relative placering i et datasæt.

\begin{definition}[Percentil]
Den $p$-te \textbf{percentil} er en værdi, som cirka $p\%$ af observationerne ligger på eller under. Der findes flere konventioner for beregning af percentiler, især for små stikprøver, så den valgte metode bør altid angives.
\end{definition}

Hvis en svartid fx ligger på 95.-percentilen, betyder det, at cirka 95\% af alle svartider ligger på eller under denne værdi.

\begin{example} Beregning af percentil

    Brug metoden med nærmeste rang til at finde 90.-percentilen for svartiderne:
\begin{enumerate}
        \item Sortér data: $\{120, 150, 180, 200, 220, 250, 280, 300, 350, 400\}$
        \item Rang: $\lceil 0.9 \times 10\rceil = 9$
        \item 90.-percentilen = 350~ms
\end{enumerate}
    \end{example}

Kvartiler er specifikke, meget anvendte percentiler, der deler data i fire lige store dele.

\begin{definition}[Kvartiler]

Kvartiler deler et ordnet datasæt i fire lige store dele.

\begin{itemize}
    \item \textbf{Q1 (første kvartil)}: 25.-percentilen.
    \item \textbf{Q2 (anden kvartil)}: 50.-percentilen. Det er datasættets \textbf{median}.
    \item \textbf{Q3 (tredje kvartil)}: 75.-percentilen.
\end{itemize}
\end{definition}

\begin{example} Beregning af kvartiler

For de sorterede svartider: $\{120, 150, 180, 200, 220, | 250, 280, 300, 350, 400\}$
\begin{itemize}
    \item Den nederste halvdel er $\{120, 150, 180, 200, 220\}$. Medianen af denne halvdel er $Q1 = 180$~ms.
    \item Den øverste halvdel er $\{250, 280, 300, 350, 400\}$. Medianen af denne halvdel er $Q3 = 300$~ms.
\end{itemize}
\end{example}

\begin{example} Kvartiler aflæst på en sumkurve

Kvartiler kan aflæses direkte på en sumkurve. Som eksempel bruger vi igen pærernes levetider. Aflæst på sumkurven i \autoref{fig:lamp_life_ogive_quartiles} er de omtrentlige kvartiler $Q1 \approx 900$, $Q2$ (medianen) $\approx 1000$ og $Q3 \approx 1100$ timer.

\begin{figure}[h!]
\centering
\begin{tikzpicture}
    \def\xQone{910.35}
    \def\xQtwo{996.55}
    \def\xQthree{1116.28}

    \begin{axis}[
        title={Kumuleret relativ hyppighed af pærernes levetider (sumkurve)},
        clip=false,
        xlabel={Levetid (i timer)},
        ylabel={Kumuleret relativ hyppighed},
        ymin=0, ymax=1.05,
        xmin=400, xmax=1500,
        width=0.9\linewidth,
        height=7cm,
        grid=major,
        grid style={dashed,gray!30},
        title style={font=\large\bfseries},
        xtick={400,600,800,1000,1200,1400},
        xticklabels={400,600,800,,{1\,200},{1\,400}},
        ytick={0.25,0.50,0.75,1.00},
        yticklabels={$\pmb{0.25}$,$\pmb{0.50}$,$\pmb{0.75}$,},
        yticklabel style={font=\bfseries, text=mseViaBlue},
        ]
        \addplot[color=mseViaBlue, thick, mark=*, mark options={fill=mseViaBlue}] coordinates {
            (500, 0)
            (600, 0.010)
            (700, 0.035)
            (800, 0.095)
            (900, 0.220)
            (1000, 0.510)
            (1100, 0.715)
            (1200, 0.930)
            (1300, 0.965)
            (1400, 0.995)
            (1500, 1.000)
        };

        \draw[mseViaBlue, dashed, thick]
        (axis cs:400,0.25) -- (axis cs:\xQone,0.25) -- (axis cs:\xQone,0);
        \draw[mseViaBlue, dashed, thick]
        (axis cs:400,0.5) -- (axis cs:\xQtwo,0.5) -- (axis cs:\xQtwo,0);
        \draw[mseViaBlue, dashed, thick]
        (axis cs:400,0.75) -- (axis cs:\xQthree,0.75) -- (axis cs:\xQthree,0);

        \node[anchor=east, font=\bfseries, text=black] at (axis cs: 400,1.00) {$1.00$};

        \node[anchor=north, font=\bfseries, text=mseViaBlue, yshift=-1.5] at (axis cs:\xQone,0) {$\pmb{Q1}$};
        \node[anchor=north, font=\bfseries, text=mseViaBlue, yshift=-1.5] at (axis cs:\xQtwo,0) {$\pmb{Q2}$};
        \node[anchor=north, font=\bfseries, text=mseViaBlue, yshift=-1.5] at (axis cs:\xQthree,0) {$\pmb{Q3}$};
    \end{axis}
\end{tikzpicture}
\caption{En sumkurve for pærernes levetider med kvartilerne $Q1$, $Q2$ og $Q3$ markeret med stiplede linjer.}
\label{fig:lamp_life_ogive_quartiles}
\end{figure}

\end{example}

Percentiler deler data i 100 lige store dele, mens kvartiler deler data i fire lige store dele.

\textbf{Kvartilafstanden} er et robust spredningsmål, som generelt er langt mindre følsomt over for ekstreme observationer end den fulde variationsbredde eller standardafvigelsen.

\begin{definition}[Kvartilafstand]
    \textbf{Kvartilafstanden (IQR)} måler spredningen af de midterste 50\% af data og er defineret som
    $$
    IQR = Q_3 - Q_1,
    $$
    hvor $Q_1$ og $Q_3$ er første og tredje kvartil. En større IQR betyder, at den centrale del af data er mere spredt, mens en mindre IQR betyder, at den er tættere samlet.
\end{definition}
    

\begin{example} Beregning af IQR

For svartiderne: $\{120, 150, 180, 200, 220, 250, 280, 300, 350, 400\}$

\begin{itemize}
    \item Nederste halvdel: $\{120, 150, 180, 200, 220\}$ → $Q1 = 180$
    \item Øverste halvdel: $\{250, 280, 300, 350, 400\}$ → $Q3 = 300$
    \item $IQR = 300 - 180 = 120$ millisekunder
\end{itemize}
\end{example}


\subsection*{Opsummerende statistik}

Før vi afslutter dette kapitel, gennemgår vi de opsummerende statistikker, som er behandlet i dette afsnit. \autoref{tab:summary_stats} er et hypotetisk eksempel for et separat datasæt med 1000 svartider fra et websted. De underliggende observationer er ikke angivet, så værdierne er medtaget for at illustrere opbygningen af en opsummerende tabel og ikke som en beregning, der kan efterprøves ud fra de ti foregående observationer.

\begin{example}

For et datasæt med 1000 svartider fra et websted:

\begin{table}[htbp]
\centering
\renewcommand{\arraystretch}{1.3}
\begin{tabular}{|l|c|}
\hline
\textbf{Mål} & \textbf{Værdi} \\
\hline
Antal & 1000 \\
Gennemsnit & 245.3 ms \\
Median & 238.0 ms \\
Typetal & 220.0 ms \\
Standardafvigelse & 95.6 ms \\
Varians & 9139.4 ms² \\
Minimum & 45.0 ms \\
Maksimum & 1200.0 ms \\
Variationsbredde & 1155.0 ms \\
Q1 & 180.0 ms \\
Q3 & 300.0 ms \\
IQR & 120.0 ms \\
Skævhed & 1.2 (positiv skævhed) \\
Kurtosis & 2.8 (leptokurtisk) \\
\hline
\end{tabular}
\caption{Fuldstændig statistisk opsummering}
\label{tab:summary_stats}
\end{table}
\end{example}

Tabellen giver et fuldstændigt billede af data og er ofte et første skridt i analysen af et datasæt. Bemærk, at skævhed gennemgås i næste afsnit, mens kurtosis ikke forklares i denne bog. Vi har valgt at medtage dem her for fuldstændighedens skyld, da de ofte bruges i analysen af et datasæt -- selv om det ligger uden for denne bogs rammer at beregne dem.


\section{Forståelse af datafordelinger} \label{sec:data_distributions}
Når vi har visualiseret vores data med histogrammer og boksplot, kan vi begynde at analysere fordelingens \textit{form}. Formen fortæller os om de underliggende mønstre i data og hjælper os med at vælge de rette statistiske værktøjer. I softwareudvikling er det afgørende at forstå fordelingen af målinger som svartider eller fejlrater, når man skal fastlægge referenceniveauer for ydelse og opdage afvigelser. Først genopfrisker vi begrebet outliers.

\subsection*{Normale datasæt}

Mange store datasæt giver histogrammer med en karakteristisk overordnet form: Hyppighederne topper nær stikprøvens median og aftager symmetrisk på begge sider, så der dannes en karakteristisk \textbf{klokkeformet kurve}. Datasæt med denne egenskab kaldes \textbf{normale}, og deres histogrammer kaldes \textbf{normale histogrammer} (se \autoref{fig:normal_hist}).

\begin{figure}[htbp]
    \centering
    \begin{tikzpicture}
    \begin{axis}[
        width=0.6\textwidth,
        height=6cm,
        ybar,
        bar width=1,
        enlargelimits=0.05,
        ylabel={Hyppighed},
        xmin=1.1,
        xmax=18,
        xtick=\empty,
        ytick=\empty,
        axis lines*=left,
        y axis line style={opacity=1},
        x axis line style={opacity=1},
        grid=major, 
        grid style={dashed,gray!30}, 
        xmajorgrids=false,
        ]
    \addplot[fill=mseViaBlue, draw=mseViaBlue!80!black] coordinates {
        (1, 1) (2, 2) (3, 4) (4, 7) (5, 9) (6, 11) (7, 13) (8, 15) (9, 17) (10, 15) (11, 13) (12, 11) (13, 9) (14, 7) (15, 4) (16, 2) (17, 1)
    };
    \end{axis}
    \end{tikzpicture}
    \caption{Et histogram over et normalt datasæt, der viser den karakteristiske klokkeform.}
    \label{fig:normal_hist}
\end{figure}


I softwaresystemer er perfekt normale fordelinger sjældne, men nogle målinger kan være \textbf{tilnærmelsesvis normale}. Svartiden (latensen) for en stærkt optimeret, intern mikrotjeneste, der udfører en ensartet opgave, kan fx være tilnærmelsesvis normal. Et histogram over et tilnærmelsesvis normalt datasæt ser ud som \autoref{fig:approx_normal_hist}.

\begin{figure}[htbp]
    \centering
    \begin{tikzpicture}
    \begin{axis}[
        width=0.6\textwidth,
        height=6cm,
        ybar,
        bar width=1,
        enlargelimits=0.05,
        ylabel={Hyppighed},
        xmin=1.1,
        xmax=18,
        xtick=\empty,
        ytick=\empty,
        axis lines*=left,
        y axis line style={opacity=1},
        x axis line style={opacity=1},
        grid=major, 
        grid style={dashed,gray!30}, 
        xmajorgrids=false,
        ]
    \addplot[fill=mseViaBlue, draw=mseViaBlue!80!black] coordinates {
        (1,1) (2,2) (3,3) (4,5) (5,7) (6,9) (7,11) (8,12) (9,10) (10,11) (11,9) (12,7) (13,4) (14,2) (15,1) (17,2)
    };
    \end{axis}
    \end{tikzpicture}
    \caption{Et histogram over et datasæt, der er tilnærmelsesvis normalt. Selv om det ikke er perfekt symmetrisk, følger det den overordnede klokkeform.}
    \label{fig:approx_normal_hist}
\end{figure}

Mange datasæt fra virkeligheden, fra menneskers højde til målefejl, følger netop denne form. Om sådanne data siger man, at de er \textbf{normalfordelte}.

\begin{definition}[Normalfordeling]
En \textbf{normalfordeling} er en symmetrisk, klokkeformet fordeling, hvor middelværdi, median og typetal alle er ens og ligger i centrum. Dens form er fuldstændig bestemt af dens middelværdi ($\mu$) og standardafvigelse ($\sigma$).
\end{definition}

Et histogram er et skøn over formen af den fordeling, som har frembragt observationerne. Med flere data og et passende valg af intervalbredde kan det give et mere detaljeret billede af den underliggende fordeling, men det bliver ikke nødvendigvis normalt. Om en normalfordelingsmodel er passende, afhænger af den proces, der frembringer data.

\begin{theorem}[Store tals lov]
    For uafhængige og identisk fordelte observationer med en endelig forventet værdi konvergerer stikprøvens gennemsnit mod populationens middelværdi, når stikprøvestørrelsen vokser.
\end{theorem}

Store tals lov handler om stikprøvens gennemsnit, ikke om et histograms form. Hvis den underliggende population selv er normalfordelt, har dens tæthedsfunktion følgende form:

\begin{figure}[htbp]
    \centering
    \begin{tikzpicture}
    \begin{axis}[
        width=0.6\textwidth,
        height=6cm,
        axis lines*=left,
        xmin=0,
        xmax=10,
        ymin=0,
        ymax=0.5,
        xtick=\empty,
        ytick=\empty,
        grid=major,
        grid style={dashed,gray!30},
        xmajorgrids=false
    ]
    % Smooth normal curve (mean at 5, sd ~1.3) in mseViaBlue
    \addplot[smooth, very thick, color=mseViaBlue, samples=200, domain=0:10] { 0.45*exp(-0.5*((x-5)/1.3)^2) };
    \end{axis}
    \end{tikzpicture}
    \caption{Normalfordelingskurven.}
    \label{fig:normal_curve_smooth}
\end{figure}

For data, der er tilnærmelsesvis normale, kan vi bruge en effektiv tommelfingerregel til at forstå spredningen uden at se på hvert eneste datapunkt. Den kaldes \textbf{den empiriske regel}.

\begin{theorem}[Den empiriske regel]
Hvis et datasæt er tilnærmelsesvis normalt med stikprøvegennemsnit $\bar{x}$ og stikprøvestandardafvigelse $s$, gælder:
\begin{itemize}
    \item Cirka \textbf{68.3\%} af observationerne ligger inden for 1 standardafvigelse fra gennemsnittet ($\bar{x} \pm s$).
    \item Cirka \textbf{95.4\%} af observationerne ligger inden for 2 standardafvigelser fra gennemsnittet ($\bar{x} \pm 2s$).
    \item Cirka \textbf{99.7\%} af observationerne ligger inden for 3 standardafvigelser fra gennemsnittet ($\bar{x} \pm 3s$).
\end{itemize}
\end{theorem}

\begin{figure}[htbp]
    \centering
    \begin{tikzpicture}
    \begin{axis}[
        width=0.6\textwidth,
        height=6cm,
        axis lines*=left,
        xmin=-3.5,
        xmax=3.5,
        ymin=0,
        ymax=0.5,
        xtick={-3,-2,-1,0,1,2,3},
        ytick=\empty,
        grid=major,
        grid style={dashed,gray!30},
        xmajorgrids=false,
        clip=false
    ]
    % Base light fill (-3σ..+3σ)
    \addplot[domain=-3:3, samples=200, draw=none, fill=mseViaBlue!15] { 0.45*exp(-0.5*(x^2)) } -- (axis cs:3,0) -- (axis cs:-3,0) -- cycle;
    % Medium fill (-2σ..+2σ)
    \addplot[domain=-2:2, samples=200, draw=none, fill=mseViaBlue!30] { 0.45*exp(-0.5*(x^2)) } -- (axis cs:2,0) -- (axis cs:-2,0) -- cycle;
    % Darker fill (-1σ..+1σ)
    \addplot[domain=-1:1, samples=200, draw=none, fill=mseViaBlue!50] { 0.45*exp(-0.5*(x^2)) } -- (axis cs:1,0) -- (axis cs:-1,0) -- cycle;
    % Normal curve outline
    \addplot[smooth, very thick, color=mseViaBlue, samples=200, domain=-3.2:3.2] { 0.45*exp(-0.5*(x^2)) };
    % Vertical sigma lines (calculated to match the curve: 0.45*exp(-0.5*x^2))
    \addplot[mseViaBlue!80!black, thick] coordinates {(-3,0) (-3,0.0050)};
    \addplot[mseViaBlue!80!black, thick] coordinates {(-2,0) (-2,0.0609)};
    \addplot[mseViaBlue!80!black, thick] coordinates {(-1,0) (-1,0.2729)};
    \addplot[mseViaBlue!80!black, thick, dashed] coordinates {(0,0) (0,0.45)};
    \addplot[mseViaBlue!80!black, thick] coordinates {(1,0) (1,0.2729)};
    \addplot[mseViaBlue!80!black, thick] coordinates {(2,0) (2,0.0609)};
    \addplot[mseViaBlue!80!black, thick] coordinates {(3,0) (3,0.0050)};
    % Percentage labels above segments
    \node[font=\scriptsize] at (axis cs:-3.2,0.03) {0.15\%};
    \node[font=\scriptsize] at (axis cs:3.2,0.03) {0.15\%};
    \node[font=\scriptsize] at (axis cs:-2.5,0.08) {2.1\%};
    \node[font=\scriptsize] at (axis cs:2.5,0.08) {2.1\%};
    \node[font=\scriptsize] at (axis cs:-1.5,0.20) {13.6\%};
    \node[font=\scriptsize] at (axis cs:1.5,0.20) {13.6\%};
    \node[font=\scriptsize] at (axis cs:-0.5,0.28) {34.15\%};
    \node[font=\scriptsize] at (axis cs:0.5,0.28) {34.15\%};
    % Braces for 68.3%, 95.4%, 99.7%
    \draw[mseViaBlue!80!black, thick, decorate, decoration={brace, amplitude=4pt, mirror}] (axis cs:-1,-0.05) -- (axis cs:1,-0.05) node[midway, below=4pt, xshift=2pt, font=\scriptsize]{68.3\%};
    \draw[mseViaBlue!80!black, thick, decorate, decoration={brace, amplitude=4pt, mirror}] (axis cs:-2,-0.11) -- (axis cs:2,-0.10) node[midway, below=4pt, xshift=2pt, font=\scriptsize]{95.4\%};
    \draw[mseViaBlue!80!black, thick, decorate, decoration={brace, amplitude=4pt, mirror}] (axis cs:-3,-0.17) -- (axis cs:3,-0.16) node[midway, below=4pt, xshift=2pt, font=\scriptsize]{99.7\%};
    \end{axis}
    \end{tikzpicture}
    \caption{Den empiriske regel vist på en normalfordelingskurve. De indre bånd markerer $\pm1\sigma$ (68.3\%), $\pm2\sigma$ (95.4\%) og $\pm3\sigma$ (99.7\%).}
    \label{fig:empirical_rule_shaded}
\end{figure}

\subsection*{Skæve datasæt}
Når et datasæt ikke er symmetrisk omkring stikprøvens median, beskrives det som \textbf{skævt}. En lang hale, der strækker sig mod højre, betyder, at fordelingen er \textbf{højreskæv}; en lang hale mod venstre betyder, at den er \textbf{venstreskæv}. Histogrammet i \autoref{fig:left_skewed_hist} er altså venstreskævt, mens histogrammet i \autoref{fig:right_skewed_hist} er højreskævt.

\begin{definition}[Skævhed]
    \textbf{Skævhed} måler asymmetrien i en datafordeling:
    \begin{itemize}
        \item \textbf{Positiv skævhed (højreskæv)}: Halen strækker sig mod højre (gennemsnit > median)
        \item \textbf{Negativ skævhed (venstreskæv)}: Halen strækker sig mod venstre (gennemsnit < median)
        \item \textbf{Symmetrisk}: Gennemsnit $\approx$ median
    \end{itemize}
    \end{definition}

\begin{figure}[htbp]
    \centering
    \begin{tikzpicture}
        \begin{axis}[
            width=0.6\textwidth,
            height=6cm,
            ybar,
            bar width=1,
            enlargelimits=0.05,
            ylabel={Hyppighed},
            xmin=1.1,
            xmax=18,
            xtick=\empty,
            ytick=\empty,
            axis lines*=left,
            y axis line style={opacity=1},
            x axis line style={opacity=1},
            grid=major, 
            grid style={dashed,gray!30}, 
            xmajorgrids=false,
            ]
        \addplot[fill=mseViaBlue, draw=mseViaBlue!80!black] coordinates {
            (1, 1) (2, 1) (4, 2) (5, 3) (6, 4) (7, 6) (8, 8) (9, 12) (10, 15) (11, 17) (12, 16) (13, 13) (14, 9) (15, 5) (16,1)
        };
        \end{axis}
    \end{tikzpicture}
    \caption{Et histogram over et datasæt, der er venstreskævt (negativt skævt). Fordelingens hale ligger til venstre.}
    \label{fig:left_skewed_hist}
\end{figure}

\begin{figure}[htbp]
    \centering
    \begin{tikzpicture}
        \begin{axis}[
            width=0.6\textwidth,
            height=6cm,
            ybar,
            bar width=1,
            enlargelimits=0.05,
            ylabel={Hyppighed},
            xmin=1.1,
            xmax=18,
            xtick=\empty,
            ytick=\empty,
            axis lines*=left,
            y axis line style={opacity=1},
            x axis line style={opacity=1},
            grid=major, 
            grid style={dashed,gray!30}, 
            xmajorgrids=false,
            ]
        \addplot[fill=mseViaBlue, draw=mseViaBlue!80!black] coordinates {
            (1, 1) (2, 5) (3, 9) (4, 13) (5, 16) (6, 17) (7, 15) (8, 12) (9, 8) (10, 6) (11, 4) (12, 3) (13, 2) (14, 1) (16, 1)
        };
        \end{axis}
    \end{tikzpicture}
    \caption{Et histogram over et datasæt, der er højreskævt (positivt skævt). Fordelingens hale ligger til højre.}
    \label{fig:right_skewed_hist}
\end{figure}

\newpage

\autoref{fig:skew_mean_median_mode} illustrerer almindelige mønstre for gennemsnit, median og typetal i unimodale skæve fordelinger. Disse rækkefølger er nyttige tommelfingerregler, ikke definitioner af skævhed og ikke generelle regler.


\begin{figure}[htbp]
    \centering
    % Positive skew panel
    \begin{minipage}{0.32\textwidth}
        \centering
        \begin{tikzpicture}
        \begin{axis}[
            width=\linewidth,
            height=4cm,
            xmin=0, xmax=10,
            ymin=0, ymax=0.5,
            axis lines*=left,
            xtick=\empty,
            ytick=\empty,
            clip=false
        ]
        % Right-tailed (positive skew) smooth curve
        \addplot[smooth, very thick, color=mseViaBlue, samples=200, domain=0:10] { (x)*exp(-0.6*x)/3 };
        % Mode, Median, Mean
        \addplot[mseViaBlue!80!black, thick, dotted] coordinates {(1.7,0) (1.7,0.46)};
        \addplot[third, thick] coordinates {(2.3,0) (2.3,0.46)};
        \addplot[mseViaBlue!80!black, thick, dashed] coordinates {(2.9,0) (2.9,0.46)};
        % Vertical labels (rotated 90 degrees left)
        \node[anchor=west, font=\scriptsize, rotate=45] at (axis cs:1,0.46) {Typetal};
        \node[anchor=west, font=\scriptsize, rotate=45] at (axis cs:2,0.46) {Median};
        \node[anchor=west, font=\scriptsize, rotate=45] at (axis cs:3,0.46) {Gennemsnit};
        \end{axis}
        \end{tikzpicture}
        \par\small Højreskæv
    \end{minipage}
    \hfill
    % Symmetric (normal) panel
    \begin{minipage}{0.32\textwidth}
        \centering
        \begin{tikzpicture}
        \begin{axis}[
            width=\linewidth,
            height=4cm,
            xmin=0, xmax=10,
            ymin=0, ymax=0.5,
            axis lines*=left,
            xtick=\empty,
            ytick=\empty,
            clip=false
        ]
        \addplot[smooth, very thick, color=mseViaBlue, samples=200, domain=0:10] { 0.45*exp(-0.5*((x-5)/1.3)^2) };
        \addplot[mseViaBlue!80!black, thick, dashed] coordinates {(5,0) (5,0.46)};
        \node[anchor=south, font=\scriptsize, align=center] at (axis cs:5,0.48) {Gennemsnit\\Median\\Typetal};
        \end{axis}
        \end{tikzpicture}
        \par\small Symmetrisk
    \end{minipage}
    \hfill
    % Negative skew panel
    \begin{minipage}{0.32\textwidth}
        \centering
        \begin{tikzpicture}
        \begin{axis}[
            width=\linewidth,
            height=4cm,
            xmin=0, xmax=10,
            ymin=0, ymax=0.5,
            axis lines*=left,
            xtick=\empty,
            ytick=\empty,
            clip=false
        ]
        % Left-tailed (negative skew) smooth curve
        \addplot[smooth, very thick, color=mseViaBlue, samples=200, domain=0:10] { ((10-x))*exp(-0.6*(10-x))/3 };
        % Mean, Median, Mode (mean left, median middle, mode right)
        \addplot[mseViaBlue!80!black, thick, dashed] coordinates {(7.1,0) (7.1,0.46)};
        \addplot[third, thick] coordinates {(7.7,0) (7.7,0.46)};
        \addplot[mseViaBlue!80!black, thick, dotted] coordinates {(8.3,0) (8.3,0.46)};
        % Vertical labels (rotated 90 degrees left)
        \node[anchor=west, font=\scriptsize, rotate=45] at (axis cs:6.5,0.46) {Gennemsnit};
        \node[anchor=west, font=\scriptsize, rotate=45] at (axis cs:7.5,0.46) {Median};
        \node[anchor=west, font=\scriptsize, rotate=45] at (axis cs:8.5,0.46) {Typetal};
        \end{axis}
        \end{tikzpicture}
        \par\small Venstreskæv
    \end{minipage}
    \caption{Sammenhængen mellem gennemsnit, median og typetal for positivt skæve, symmetriske og negativt skæve fordelinger.}
    \label{fig:skew_mean_median_mode}
\end{figure}

\subsection*{Identifikation af outliers}

Som nævnt i \autoref{sec:ch6sec1} er outliers datapunkter, der adskiller sig markant fra resten af data. De kan opstå på grund af målefejl, indtastningsfejl eller reel variation i data. Det er vigtigt at finde outliers, fordi de kan forvrænge statistiske analyser.

En almindelig metode til at identificere outliers bruger kvartilafstanden (IQR):

\begin{itemize}
    \item \textbf{Moderate outliers}: Datapunkter, der ligger uden for $1.5\times IQR$-grænserne, men ikke uden for $3\times IQR$-grænserne.
    \item \textbf{Ekstreme outliers}: Datapunkter, der er mindre end $Q_1 - 3 \times IQR$ eller større end $Q_3 + 3 \times IQR$.
\end{itemize}

\begin{example} Identifikation af outliers

Givet svartiderne: $\{120, 150, 180, 200, 220, 250, 280, 300, 350, 400\}$

\begin{itemize}
    \item $Q_1 = 180$, $Q_3 = 300$, $IQR = 120$
    \item Nedre grænse: $Q_1 - 1.5 \times IQR = 180 - 180 = 0$
    \item Øvre grænse: $Q_3 + 1.5 \times IQR = 300 + 180 = 480$
\end{itemize}

Alle datapunkter ligger mellem 0 og 480, så der er ingen moderate outliers i dette datasæt.

\end{example}

\newpage

\begin{example}

Antag, at vi har følgende svartider (i millisekunder) for en webapplikation:

$\{120, 150, 180, 200, 220, 250, 280, 300, 350, 400, 1200\}$.

Lad os undersøge, om der er en øvre outlier:

\begin{itemize}
    \item Sorterede data: $120, 150, 180, 200, 220, 250, 280, 300, 350, 400, 1200$
    \item $n = 11$
    \item Median ($Q_2$): 250 (6. værdi)
    \item Nederste halvdel: $120, 150, 180, 200, 220$ $\rightarrow Q_1 = 180$
    \item Øverste halvdel: $280, 300, 350, 400, 1200$ $\rightarrow Q_3 = 350$
    \item $IQR = Q_3 - Q_1 = 350 - 180 = 170$
    \item Øvre grænse: $Q_3 + 1.5 \times IQR = 350 + 1.5 \times 170 = 350 + 255 = 605$
\end{itemize}

Værdien $1200$ er større end $605$, så den er en øvre outlier.

\end{example}

\subsection*{Boksplot}

Boksplot giver en kompakt opsummering af en datafordeling og fremhæver medianen, kvartilerne og outliers. \autoref{fig:boxplot_from_image} viser et eksempel på et boksplot.

\begin{figure}[htbp]
    \centering
    \resizebox{0.7\textwidth}{!}{% Danish redrawing of the annotated box plot (same labels as figure/boxplot.png).
\begin{tikzpicture}[x=1cm, y=1cm, font=\small,
    lbl/.style={draw=black!40, fill=black!4, align=center, inner sep=4pt}]
  \definecolor{bxorange}{RGB}{217,108,40}
  % fences (1.5 x IQR) and whiskers
  \draw[bxorange, thick, densely dotted] (-3.3,-0.6) -- (-3.3,0.6);
  \draw[bxorange, thick, densely dotted] (8.3,-0.6) -- (8.3,0.6);
  \draw[bxorange, thick] (-2.5,0) -- (1.5,0) (4.5,0) -- (7.5,0);
  \draw[bxorange, thick] (-2.5,-0.5) -- (-2.5,0.5) (7.5,-0.5) -- (7.5,0.5);
  % box and median
  \draw[bxorange, thick, fill=bxorange!12] (1.5,-1) rectangle (4.5,1);
  \draw[bxorange, very thick] (3.3,-1) -- (3.3,1);
  \node[font=\footnotesize] at (3.3,-1.35) {Median};
  \node[anchor=south] at (1.5,1.05) {$Q_1$};
  \node[anchor=south] at (4.5,1.05) {$Q_3$};
  % IQR bracket
  \draw[thick] (1.5,2.3) -- (4.5,2.3) (1.5,2.1) -- (1.5,2.5) (4.5,2.1) -- (4.5,2.5);
  \node[font=\bfseries] at (3,2.95) {Kvartilafstand (IQR)};
  % outliers
  \foreach \x in {-5.6,-4.9,-4.2} \draw[bxorange, thick] (\x,0) circle (3pt);
  \foreach \x in {9.4,10.1,11.0} \draw[bxorange, thick] (\x,0) circle (3pt);
  \draw (-5.75,-0.35) to[out=-90,in=90] (-5.75,-0.45) -- (-4.05,-0.45) to[out=90,in=-90] (-4.05,-0.35);
  \draw (9.25,-0.35) to[out=-90,in=90] (9.25,-0.45) -- (11.15,-0.45) to[out=90,in=-90] (11.15,-0.35);
  \node at (-4.9,-0.9) {outliers};
  \node at (10.2,-0.9) {outliers};
  % labelled boxes
  \node[lbl] (lo) at (-3.3,2.3) {Nedre $1.5\times\mathrm{IQR}$-knurhår\\$(Q_1 - 1.5 \times \mathrm{IQR})$};
  \node[lbl] (hi) at (8.3,2.3) {Øvre $1.5\times\mathrm{IQR}$-knurhår\\$(Q_3 + 1.5 \times \mathrm{IQR})$};
  \draw[black!40] (lo.south) -- (-3.3,0.7) (hi.south) -- (8.3,0.7);
  \node[lbl] (mn) at (-1.4,-2.3) {Mindste værdi};
  \node[lbl] (mx) at (6.6,-2.3) {Største værdi};
  \draw[black!40] (mn.north) -- (-2.5,-0.05) (mx.north) -- (7.5,-0.05);
\end{tikzpicture}
}
    \caption{Et eksempel på et boksplot, der illustrerer medianen, kvartilerne, knurhårene og mulige outliers.}
    \label{fig:boxplot_from_image}
\end{figure}


\begin{definition}[Et boksplots bestanddele]
Et boksplot bygger på femtalsopsummeringen og kan også vise outliers:
\begin{itemize}
    \item \textbf{Nedre knurhår}: Som regel den mindste observation inden for den nedre $1.5\times IQR$-grænse
    \item \textbf{Q1 (første kvartil)}: 25.-percentilen
    \item \textbf{Median (Q2)}: 50.-percentilen
    \item \textbf{Q3 (tredje kvartil)}: 75.-percentilen
    \item \textbf{Øvre knurhår}: Som regel den største observation inden for den øvre $1.5\times IQR$-grænse
    \item \textbf{Outliers}: Punkter, der ligger mere end 1.5 × IQR fra boksen
\end{itemize}
\end{definition}

Et boksplot giver en visuel opsummering af et datasæts fordeling, så det er let at sammenligne grupper og få øje på outliers. Sådan læser man et boksplot:

\begin{itemize}
    \item \textbf{Boksens længde} viser spredningen af den midterste halvdel af data. En længere boks betyder mere variation i de midterste 50\%.
    \item \textbf{Medianlinjens placering} og knurhårenes relative længder kan antyde asymmetri, men et boksplot alene fastlægger ikke retningen eller graden af skævhed.
    \item I mange unimodale skæve fordelinger trækkes gennemsnittet mod den længste hale. Derfor ser man ofte $\text{typetal}<\text{median}<\text{gennemsnit}$ for højreskæve data og den omvendte rækkefølge for venstreskæve data. Det er almindelige mønstre, ikke garanterede identiteter.
    \item \textbf{Knurhårene} viser variationsbredden for hovedparten af data, når outliers er udeladt.
    \item \textbf{Outliers} fremhæver usædvanligt høje eller lave værdier, som kan kræve nærmere undersøgelse.
\end{itemize}


Boksplot er særligt nyttige til at sammenligne fordelinger på tværs af flere grupper eller datasæt, da de hurtigt afslører forskelle i medianer, spredning og forekomst af outliers.

Dette afslutter vores gennemgang af datavisualisering og opsummerende statistik og dermed vores gennemgang af sandsynlighedsregning og statistik i det hele taget.
